\begin{abstract}
La presente tesi si colloca nell'ambito dello sviluppo di intelligenze artificiali per la creazione di \textbf{Non-Player Characters} (\textbf{NPC})\cite{nextgen1996}, caratterizzati da personalità realistiche e complesse. In particolare, viene affrontata la carenza di strumenti in grado di generare dialoghi realistici e coerenti in tempo reale, senza l'intervento umano.\\\\
Più in dettaglio, il lavoro si concentra sull'uso di \textbf{Large Language Models} (\textbf{LLM})\cite{brown1993mathematics} per la simulazione del comportamento, seguendo il modello dei \textbf{Big Five} di \textbf{Lewis Goldberg} (\textbf{GBF})\cite{goldberg1992} per la definizione e valutazione dei tratti di personalità.
L'obiettivo della tesi è sviluppare un'applicazione in grado di generare NPC (Non-Player-Characters) con personalità coerenti e realistiche, dotati di tratti di personalità basati sul GBF e su una \textbf{storia di background personalizzata}.\\
Il sistema è stato progettato per essere compatibile con qualsiasi LLM, a patto che esso supporti una \textbf{Context Length} (\textbf{CL}) superiore a \textbf{4000} token. Questo requisito risulta essenziale poiché i prompt e le risposte generati mediamente superano i 3800 token. Una CL adeguata assicura la flessibilità nell'integrazione di diverse architetture di modelli, preservando la coerenza e la qualità delle interazioni.\\\\
\textbf{I contributi principali di questa tesi includono:}
\begin{enumerate}
    \item Un sistema di generazione, sia manuale che automatico, di NPC con tratti di personalità coerenti;
    \item La valutazione dell'accuratezza nella replicazione dei tratti psicologici mediante un test dell'\textbf{International Personality Item Pool} (\textbf{IPIP})\cite{goldberg2006} per i GBF;
    \item Un confronto dettagliato tra 3 diversi modelli con diversi livelli e tecniche di quantizzazione in termini di precisione nella replicazione e velocità di inferenza.
\end{enumerate}

\textbf{Durante lo sviluppo le sfide tecniche sono state:}
\begin{enumerate}
    \item L'Applicazione di un test psicometrico sugli LLM;
    \item Il Bilanciamento tra tecniche di prompting efficaci e tempi di inferenza.
\end{enumerate}
\end{abstract} 